\documentclass[11pt,a4paper]{article}

\usepackage[utf8]{inputenc}
\usepackage[T1]{fontenc}
\usepackage{lmodern}
\usepackage{amsmath,amssymb,amsthm}
\usepackage{booktabs}
\usepackage{array}
\usepackage{caption}
\usepackage{geometry}
\usepackage{hyperref}
\usepackage{listings}
\usepackage{xcolor}
\usepackage{microtype}

\geometry{margin=1in}
\setlength{\emergencystretch}{3em}

\captionsetup[lstlisting]{singlelinecheck=false,justification=raggedright}

\hypersetup{
  colorlinks=true,
  linkcolor=blue!60!black,
  citecolor=green!50!black,
  urlcolor=blue!60!black
}

\lstdefinelanguage{Lean4}{
  morekeywords={def,theorem,lemma,abbrev,by,rfl,let,in,if,then,else,match,with,fun,true,false,set_option},
  sensitive=true,
  morecomment=[l]{--},
  morecomment=[s]{/-}{-/},
  morestring=[b]"
}

\lstset{
  language=Lean4,
  basicstyle=\ttfamily\footnotesize,
  keywordstyle=\color{blue!70!black}\bfseries,
  commentstyle=\color{gray!70!black}\itshape,
  stringstyle=\color{green!50!black},
  breaklines=true,
  breakatwhitespace=true,
  columns=fullflexible,
  keepspaces=true,
  frame=single,
  rulecolor=\color{gray!35},
  backgroundcolor=\color{gray!5},
  xleftmargin=0.5em,
  xrightmargin=0.5em,
  aboveskip=0.8em,
  belowskip=0.8em,
  literate={⟨}{{$\langle$}}1 {⟩}{{$\rangle$}}1 {×}{{$\times$}}1 {→}{{$\to$}}1 {∧}{{$\land$}}1 {·}{{$\cdot$}}1 {±}{{$\pm$}}1
}

\newtheorem{theorem}{Theorem}[section]
\newtheorem{proposition}[theorem]{Proposition}
\newtheorem{definition}[theorem]{Definition}
\newtheorem{remark}[theorem]{Remark}

\title{\textbf{Constructive Formalization of \texorpdfstring{$\operatorname{Spin}(8)$}{Spin(8)} Triality, \texorpdfstring{$W(D_4)$}{W(D4)} Chiral Orthogonality, and Mumford Spinor Representations in Lean~4}}

\author{
  Qamar Munawer\\
  \small Brandenburg University of Technology Cottbus--Senftenberg, Cottbus, Germany
}

\date{}

\begin{document}
\maketitle

\begin{abstract}
This paper presents a constructive, axiom-free formalization in Lean~4 of the finite-dimensional Lie-theoretic weight combinatorics, Weyl group representation theory, and root-operator linear algebra underlying Mumford's spinor Hodge classes on abelian varieties of $D_4$ type. The general Hodge Conjecture remains open, and this work formalizes strictly the discrete weight spaces, exterior-power chiral splittings, and exact group-character orthogonality of $\operatorname{Spin}(8)$ representations. With no external library imports, the Lean~4 kernel generates the $8$-dimensional vector and half-spinor weight spaces from first principles, counts the $32$ triality zero-sum triples, verifies the Hodge-star splitting of the $70$-dimensional fourth exterior power into $35_+ \oplus 35_-$ and the invertibility of the $6 \times 6$ Weyl-to-Weil transition matrix, certifies full rank $6$ for the $D_4$ simple-root Gram matrix on the Cartan zero-weight subspace, and evaluates exact character inner products over all $192$ elements of $W(D_4)$.
\end{abstract}

\section{Introduction and Epistemic Scope}

Abelian fourfolds of Mumford type~\cite{Mumford1969,Weil1977} provide the canonical examples of exceptional Hodge classes that are not generated by intersections of divisor classes. For a general abelian fourfold whose Mumford--Tate group is isogenous to $\operatorname{Spin}(8)$ (or a $\mathbb{Q}$-form of $D_4$), the first cohomology $H^1(A, \mathbb{Q})$ carries an $8$-dimensional spinor representation, and the middle cohomology $H^4(A, \mathbb{Q}) \cong \bigwedge^4 H^1(A, \mathbb{Q})$ decomposes under $\mathfrak{so}(8)$ into irreducible representations including the $35$-dimensional self-dual and anti-self-dual components $35_+ \oplus 35_-$. Weil's Hodge classes~\cite{Weil1977,vanGeemen1994,Zarhin1984} arise inside the $6$-dimensional Cartan zero-weight space $(\bigwedge^4 S^+_8)_0$, where triality relates the vector representation $V_8$ and the two half-spinor representations $S^+_8$ and $S^-_8$~\cite{Voisin2002}.

A systematic audit of \texttt{mathlib4} and existing Lean~4 formalizations shows zero prior implementations of $\operatorname{Spin}(8)$ triality, the $192$-element Weyl group $W(D_4) \cong S_4 \ltimes (\mathbb{Z}/2\mathbb{Z})^3$, Hodge-star chiral splitting on $\bigwedge^4 S^+_8$, or Jordan--Wigner root-action Gram matrices. This paper documents a constructive, axiom-free formalization of these finite-dimensional Lie-theoretic invariants in Lean~4.35.0 across \texttt{ShiodaFormalConjectures.lean} and \texttt{HodgeVerification.lean}.

\begin{remark}[Epistemic Boundary]
  The formal theorems verified in Lean~4 establish finite-dimensional representation-theoretic identities, weight-space cardinalities, matrix invertibility, and Weyl character inner products for $\mathfrak{so}(8) = D_4$. They do not claim a proof of the general Hodge Conjecture~\cite{Deligne2000} or of the Mumford--Weil Hodge Conjecture in arbitrary families.
\end{remark}

\section{Constructive Generation of \texorpdfstring{$\operatorname{Spin}(8)$}{Spin(8)} Weight Spaces and Triality Zero-Sum Triples}

Using scaled integer coordinates ($2\times$ the standard half-integer lattice $\frac{1}{2}(\pm 1, \pm 1, \pm 1, \pm 1)$), the $16$ sign tuples in $\{\pm 1\}^4$ are generated algorithmically by \texttt{allSign4Tuples}. Filtering by sign product $+1$ yields the $8$ weights of the even half-spinor representation $S^+_8$ (\texttt{so8Weights}), while sign product $-1$ yields the $8$ weights of the odd half-spinor representation $S^-_8$ (\texttt{so8OddSpinorWeights}). The $8$ weights of the vector representation $V_8$ (\texttt{so8VectorWeights}) are the coordinate vectors $\pm 2 e_i$ for $0 \le i \le 3$.

The even Clifford algebra $C^+(V_8)$ has dimension $\sum_{k=0}^4 \binom{8}{2k} = 128$, matching $\dim \operatorname{End}(S^+_8) + \dim \operatorname{End}(S^-_8) = 8^2 + 8^2 = 128$ (\texttt{clifford\_even\_dim\_eq}). Furthermore, searching all $8^3 = 512$ triples in $V_8 \times S^+_8 \times S^-_8$ for weight-sum zero ($v + s_+ + s_- = 0$) isolates exactly $32$ triality-invariant zero-sum triples (\texttt{d4\_triality\_zero\_triples\_len}), reflecting the non-degenerate trilinear triality pairing $V_8 \otimes S^+_8 \otimes S^-_8 \to \mathbb{C}$.

\begin{lstlisting}[caption={Constructive $\operatorname{Spin}(8)$ weight spaces in \texttt{ShiodaFormalConjectures.lean} and triality triples in \texttt{HodgeVerification.lean}.}]
def allSign4Tuples : List Weight4 :=
  [1, -1].foldl (fun acc0 s0 =>
    [1, -1].foldl (fun acc1 s1 =>
      [1, -1].foldl (fun acc2 s2 =>
        [1, -1].foldl (fun acc3 s3 =>
          acc3 ++ [(s0, s1, s2, s3)]
        ) acc2
      ) acc1
    ) acc0
  ) []

/-- Even half-spinor weights S^+_8 (product of signs = +1). -/
def so8Weights : List Weight4 :=
  allSign4Tuples.filter (fun w => intEq (w.1 * w.2.1 * w.2.2.1 * w.2.2.2) 1)

/-- Odd half-spinor weights S^-_8 (product of signs = -1). -/
def so8OddSpinorWeights : List Weight4 :=
  allSign4Tuples.filter (fun w => intEq (w.1 * w.2.1 * w.2.2.1 * w.2.2.2) (-1))

/-- Vector representation weights V_8 (± 2 e_i). -/
def so8VectorWeights : List Weight4 :=
  [0, 1, 2, 3].foldl (fun acc i =>
    acc ++ [2, -2].map (fun s => basisVec4 i s)
  ) []

def clifford_even_dim : Nat := binom 8 0 + binom 8 2 + binom 8 4 + binom 8 6 + binom 8 8
def chevalley_spinor_end_dim : Nat := 8 * 8 + 8 * 8
theorem clifford_even_dim_eq : (clifford_even_dim == 128 && chevalley_spinor_end_dim == 128) = true := by rfl

def triality_all_triples : List (Weight × Weight × Weight) :=
  so8_vector_weights.foldr (fun v acc1 =>
    so8_weights.foldr (fun sp acc2 =>
      (so8_odd_spinor_weights.map (fun sm => (v, sp, sm))) ++ acc2
    ) acc1
  ) []

def triality_zero_triples : List (Weight × Weight × Weight) :=
  triality_all_triples.filter (fun t => is_zero_weight (add_w t.1 (add_w t.2.1 t.2.2)))

theorem d4_triality_zero_triples_len : triality_zero_triples.length = 32 := by rfl
\end{lstlisting}

\section{Exterior Power \texorpdfstring{$\bigwedge^4 S^+_8 \cong 35_+ \oplus 35_-$}{Wedge4(S+8) = 35+ + 35-}, Cartan Zero-Weight States, and the \texorpdfstring{$W(D_4)$}{W(D4)}-to-Weil Transition Matrix}

The fourth exterior power $\bigwedge^4 S^+_8$ has dimension $\binom{8}{4} = 70$ (\texttt{four\_tuples\_len}). Filtering the $70$ four-element subsets of $S^+_8$ for zero weight sum produces the $6$-dimensional Cartan zero-weight subspace $(\bigwedge^4 S^+_8)_0$ (\texttt{zero\_tuples\_len}). For any $4$-element subset $J \subset \{0, \dots, 7\}$, its complement $J^c = \texttt{complement8 } J$ satisfies $(J^c)^c = J$ and $J^c \neq J$, and the shuffle permutation sign $\operatorname{sgn}(J \,\|\, J^c)$ computed by \texttt{hodgeStarSign} satisfies $\star^2 = +1$ identically across all $70$ subsets. Pairing each $J$ with $J^c$ partitions the $70$ basis elements into $35$ self-dual ($+$) and $35$ anti-self-dual ($-$) pairs, and partitions the $6$ zero-weight states into $3$ complementary pairs \texttt{zeroStarPairs}.

From the $3$ pairs $(J_k, J_k^c)$ in \texttt{zeroStarPairs}, \texttt{bWeyl} constructs the $6 \times 6$ Weyl-to-Weil transition matrix $B_{\mathrm{Weyl}}$ whose $(2k)$ and $(2k+1)$ rows are $e_{J_k} + e_{J_k^c}$ and $e_{J_k} - e_{J_k^c}$:
\begin{equation}
  B_{\mathrm{Weyl}} = \begin{pmatrix}
    e_{J_0} + e_{J_0^c} \\ e_{J_0} - e_{J_0^c} \\
    e_{J_1} + e_{J_1^c} \\ e_{J_1} - e_{J_1^c} \\
    e_{J_2} + e_{J_2^c} \\ e_{J_2} - e_{J_2^c}
  \end{pmatrix}.
\end{equation}
Since each pair $(J_k, J_k^c)$ is disjoint and the $3$ pairs partition the $6$ zero-weight basis vectors, the rows of $B_{\mathrm{Weyl}}$ are pairwise orthogonal of squared Euclidean norm $2$. Setting $B_{\mathrm{adj}} = -4 B_{\mathrm{Weyl}}$, theorem \texttt{weyl\_d4\_transition\_matrix\_invertible\_cert} proves $B_{\mathrm{Weyl}} B_{\mathrm{adj}} = -8 I_6$, establishing exact invertibility over $\mathbb{Q}$. Moreover, the sum of the three self-dual rows defines the complete-intersection polarization $\Theta^2$, while the primitive Weil classes $w_1$ and $w_2$ are orthogonal to $\Theta^2$ (\texttt{weil\_w1\_orthogonal\_to\_theta2}, \texttt{weil\_w2\_orthogonal\_to\_theta2}). Branching from $\mathfrak{so}(8)$ to $\mathfrak{su}(2,2)$ splits the $8$ zero-weight tuples of $D_3$ into $2$ pure Weil tuples and $6$ mixed $D_4$ tuples (\texttt{spinor\_to\_weil\_su22\_branching\_eq}), and the $16$ conjugate Weil axes decompose evenly across the $3$ coordinate pairs (\texttt{weyl\_d4\_conjugate\_weil\_axes\_eq}).

\begin{lstlisting}[caption={Cartan zero-weight states, Hodge-star pairs, $B_{\mathrm{Weyl}}$ invertibility, and Weil orthogonality.}]
def four_tuples : List (List Weight) := combs 4 so8_weights
def zero_weight_tuples : List (List Weight) :=
  four_tuples.filter (fun t => is_zero_weight (tuple_sum t))

theorem four_tuples_len : four_tuples.length = 70 := by rfl
theorem zero_tuples_len : zero_weight_tuples.length = 6 := by rfl

def complement8 (I : List Nat) : List Nat :=
  (rangeList 0 8).filter (fun x => !I.contains x)

def countGreater (x : Nat) : List Nat → Nat
  | [] => 0
  | y :: ys => (if x > y then 1 else 0) + countGreater x ys

def inversionCount : List Nat → Nat
  | [] => 0
  | x :: xs => countGreater x xs + inversionCount xs

def hodgeStarSign (I : List Nat) : Int :=
  if inversionCount (I ++ complement8 I) % 2 == 0 then 1 else -1

def zeroStarPairs : List (List Nat × List Nat) :=
  zeroIndexTuples.foldl (fun acc I =>
    let c := complement8 I
    if lexLtNatList I c then acc ++ [(I, c)] else acc
  ) []

def bWeyl : List (List Int) :=
  zeroStarPairs.foldl (fun acc (I, c) =>
    let i := indexOfList I zeroIndexTuples
    let j := indexOfList c zeroIndexTuples
    acc ++ [basisRow6 i j 1 1, basisRow6 i j 1 (-1)]
  ) []

def bAdjCols : List (List Int) := scalarMulMat (-4) bWeyl
def neg8_I6 : List (List Int) := diagMatrix 6 (-8)

theorem weyl_d4_transition_matrix_invertible_cert :
    zeroStarPairs.length = 3 ∧
    matMul bWeyl bAdjCols = neg8_I6 :=
  ⟨by rfl, by rfl⟩

theorem weil_w1_orthogonal_to_theta2 : dot weil_w1 theta2 = 0 := by rfl
theorem weil_w2_orthogonal_to_theta2 : dot weil_w2 theta2 = 0 := by rfl
\end{lstlisting}

\begin{lstlisting}[caption={$\mathfrak{su}(2,2)$ branching and conjugate Weil axes theorems in \texttt{HodgeVerification.lean}.}]
theorem spinor_to_weil_su22_branching_eq :
  (d3_zero_weight_tuples.length == 8 &&
   pure_weil_tuples.length == 2 &&
   mixed_d4_zero_tuples.length == 6 &&
   mixed_d4_zero_tuples.length == zero_weight_tuples.length &&
   pure_weil_tuples.length + mixed_d4_zero_tuples.length == d3_zero_weight_tuples.length) = true := by rfl

theorem weyl_d4_conjugate_weil_axes_eq :
  (coord_pure_weil_tuples.length == 8 &&
   hadamard_pure_weil_tuples.length == 8 &&
   coord_pure_weil_tuples.length + hadamard_pure_weil_tuples.length == 16 &&
   binom 8 2 == 28 &&
   binom 6 2 + 1 == 16 &&
   shimura_coset_dim == 12 &&
   k22_pair0_count == 2 &&
   k22_pair1_count == 2 &&
   k22_pair2_count == 2 &&
   k22_pair0_count + k22_pair1_count + k22_pair2_count == zero_weight_tuples.length) = true := by rfl
\end{lstlisting}

\section{The \texorpdfstring{$D_4$}{D4} Simple-Root Gram Matrix and Absence of Trivial \texorpdfstring{$\mathfrak{so}(8)$}{so(8)} Singlets}

To determine whether the $6$-dimensional Cartan zero-weight space $(\bigwedge^4 S^+_8)_0$ contains any trivial $\mathfrak{so}(8)$-invariant singlet, it suffices to compute the joint kernel of the four simple-root raising operators $E_\alpha$ for $\alpha \in \Delta(D_4)$. In scaled coordinates, the four simple roots of $D_4$ are
\begin{equation}
  \alpha_1 = 2e_0 - 2e_1, \quad \alpha_2 = 2e_1 - 2e_2, \quad \alpha_3 = 2e_2 - 2e_3, \quad \alpha_4 = 2e_2 + 2e_3,
\end{equation}
encoded by \texttt{d4SimpleRoots}. The Jordan--Wigner fermionic action \texttt{applyRootToTuple} shifts each weight $w_k$ ($k \in J$) to $w_k + \alpha$, drops terms that leave $S^+_8$ or collide with an existing index in $J$, sorts the resulting $4$-tuple, and records the exact exterior-algebra sign $(-1)^{\operatorname{inv}(\text{newLst})}$.

Stacking the actions of all four simple roots into the $6 \times 6$ Gram matrix
\begin{equation}
  G_{ab} = \sum_{\alpha \in \Delta(D_4)} \langle E_\alpha(e_{J_a}), E_\alpha(e_{J_b}) \rangle
\end{equation}
via \texttt{d4\_root\_gram\_matrix} and applying fraction-free integer Gaussian elimination yields full rank $6$, certified by \texttt{mumford\_stacked\_root\_rank == 6}. Because a $6 \times 6$ Gram matrix $M^T M$ has full rank $6$, the stacked root operator $M : (\bigwedge^4 S^+_8)_0 \to \bigoplus_{\alpha \in \Delta} \bigwedge^4 S^+_8$ is injective:
\begin{equation}
  \bigcap_{\alpha \in \Delta(D_4)} \ker E_\alpha = \{0\} \quad \text{on } \left(\bigwedge^4 S^+_8\right)_0.
\end{equation}
Hence no non-zero $\mathfrak{so}(8)$-singlet exists in $\bigwedge^4 S^+_8$; the $6$-dimensional zero-weight space is spanned entirely by non-trivial $3_+ \oplus 3_-$ Cartan subspaces inside $35_+ \oplus 35_-$.

\begin{lstlisting}[caption={Simple-root action, $D_4$ Gram matrix, and full-rank singlet exclusion in \texttt{HodgeVerification.lean}.}]
def d4SimpleRoots : List Weight := [
  addW4 (basisVec4 0 2) (basisVec4 1 (-2)),
  addW4 (basisVec4 1 2) (basisVec4 2 (-2)),
  addW4 (basisVec4 2 2) (basisVec4 3 (-2)),
  addW4 (basisVec4 2 2) (basisVec4 3 2)
]

def applyRootToTuple (alpha : Weight) (I : List Nat) : List (List Nat × Int) :=
  [0, 1, 2, 3].foldl (fun acc pos =>
    let k := getNatAt I pos
    let w2 := addW4 (getWAt so8Weights k) alpha
    match indexOfWeight w2 so8Weights with
    | none => acc
    | some k2 =>
      if I.contains k2 then acc
      else
        let newLst := replaceAtNat I pos k2
        let sgn : Int := if inversionCount newLst % 2 == 0 then 1 else -1
        acc ++ [(sortNatList newLst, sgn)]
  ) []

def d4_root_gram_matrix : List (List Int) :=
  let imgs := d4SimpleRoots.map (fun alpha => zeroIndexTuples.map (applyRootToTuple alpha))
  (rangeList 0 6).map (fun a =>
    (rangeList 0 6).map (fun b =>
      imgs.foldl (fun acc rootImgs =>
        let ta := match rootImgs.drop a with | x :: _ => x | [] => []
        let tb := match rootImgs.drop b with | x :: _ => x | [] => []
        acc + dotTupleTerms ta tb
      ) 0
    )
  )

def mumford_stacked_root_rank : Nat :=
  intMatRankAuxClean 6 d4_root_gram_matrix

set_option maxRecDepth 4096 in
theorem weyl_d4_spinor_weight_decomposition_eq :
  (four_tuples.length == 70 &&
   fourIndexTuples.all (fun I =>
     complement8 (complement8 I) == I &&
     complement8 I != I &&
     intEq (hodgeStarSign (complement8 I)) (hodgeStarSign I)) &&
   wedge4_spinor_self_dual_dim == 35 &&
   wedge4_spinor_anti_self_dual_dim == 35 &&
   wedge4_spinor_self_dual_dim + wedge4_spinor_anti_self_dual_dim == four_tuples.length &&
   zero_weight_tuples.length == 6 &&
   mumford_stacked_root_rank == zero_weight_tuples.length &&
   weyl_d4_group_order == 192 &&
   triality_zero_triples.length == 32) = true := by rfl
\end{lstlisting}

\section{Full 192-Element Weyl Group \texorpdfstring{$W(D_4)$}{W(D4)} Generation and Chiral Character Orthogonality}

The Weyl group of type $D_4$ is the semidirect product $W(D_4) = S_4 \ltimes (\mathbb{Z}/2\mathbb{Z})^3$ of order $4! \cdot 2^3 = 192$. In \texttt{HodgeVerification.lean}, \texttt{weyl\_d4\_elements} generates all $192$ group elements explicitly by pairing the $24$ permutations of $\{0, 1, 2, 3\}$ (\texttt{all\_s4\_perms}) with the $8$ even-sign-change vectors in \texttt{so8Weights}. For each pair $(p, s) \in W(D_4)$, \texttt{wd4Perm8} computes the induced permutation of the $8$ spinor weights $S^+_8$, and \texttt{wd4\_all\_even\_parity} verifies that every element of $W(D_4)$ acts as an even permutation on $S^+_8$ (i.e., $W(D_4) \hookrightarrow A_8 \subset S_8$).

Because every $g \in W(D_4)$ preserves $S^+_8$ with even parity, it commutes with the Hodge star $\star$ on $\bigwedge^4 S^+_8$ and preserves the chiral decomposition $(\bigwedge^4 S^+_8)_0 = V_+ \oplus V_-$ into $3$-dimensional self-dual ($V_+$) and anti-self-dual ($V_-$) eigenspaces. Evaluating \texttt{wd4ChiralTraces p s} across all $192$ elements computes the exact character inner products
\begin{equation}
  \sum_{g \in W(D_4)} \chi_+(g)^2 = 192, \qquad \sum_{g \in W(D_4)} \chi_-(g)^2 = 192, \qquad \sum_{g \in W(D_4)} \chi_+(g) \chi_-(g) = 0,
\end{equation}
proving by Schur's character orthogonality relations that $V_+$ and $V_-$ are non-isomorphic, irreducible $3$-dimensional representations of $W(D_4)$. Table~\ref{tab:wd4_summary} summarizes the verified Lie-theoretic and Weyl-group invariants.

\begin{table}[htbp]
  \centering
  \small
  \setlength{\tabcolsep}{3pt}
  \caption{Lean~4 kernel-verified $\operatorname{Spin}(8)$ and $W(D_4)$ invariants on $\bigwedge^4 S^+_8$.}
  \label{tab:wd4_summary}
  \begin{tabular}{lcc}
    \toprule
    Invariant / Representation & Dimension / Value & Lean~4 Certificate \\
    \midrule
    Even Clifford algebra $C^+(V_8)$ & $128 = 8^2 + 8^2$ & \footnotesize\texttt{clifford\_even\_dim\_eq} \\
    Triality zero-sum triples & $32$ & \footnotesize\texttt{d4\_triality\_zero\_triples\_len} \\
    Exterior power $\bigwedge^4 S^+_8$ & $70 = 35_+ \oplus 35_-$ & \footnotesize\texttt{weyl\_d4\_spinor\_weight\_decomposition\_eq} \\
    Cartan zero-weight subspace & $6 = 3_+ \oplus 3_-$ & \footnotesize\texttt{zero\_tuples\_len} \\
    $D_4$ simple-root Gram rank & $6$ (no singlet) & \footnotesize\texttt{weyl\_d4\_spinor\_weight\_decomposition\_eq} \\
    Weyl group $W(D_4)$ order & $192$ (even on $S^+_8$) & \footnotesize\texttt{weyl\_d4\_even\_parity\_and\_chiral\_orthogonality\_eq} \\
    Character norms $(\chi_+, \chi_-)$ & $192,\; 192,\; 0$ & \footnotesize\texttt{weyl\_d4\_even\_parity\_and\_chiral\_orthogonality\_eq} \\
    \bottomrule
  \end{tabular}
\end{table}

\begin{lstlisting}[caption={Full $192$-element $W(D_4)$ generation, even parity, and chiral character orthogonality.}]
def all_s4_perms : List (List Nat) := permsNat [0, 1, 2, 3]

def weyl_d4_elements : List (List Nat × Weight) :=
  all_s4_perms.foldl (fun acc p =>
    acc ++ so8Weights.map (fun s => (p, s))
  ) []

def wd4_all_even_parity : Bool :=
  weyl_d4_elements.all (fun (p, s) => inversionCount (wd4Perm8 p s) % 2 == 0)

def wd4ChiralTraces (p : List Nat) (s : Weight) : Int × Int :=
  let perm8 := wd4Perm8 p s
  zeroStarPairs.foldl (fun (tp, tm) (I, comp) =>
    let hs := hodgeStarSign I
    let img := I.map (getNatAt perm8)
    let sgnSort : Int := if inversionCount img % 2 == 0 then 1 else -1
    let sortedImg := sortNatList img
    if sortedImg == I then
      (tp + sgnSort, tm + sgnSort)
    else if sortedImg == comp then
      let imgC := comp.map (getNatAt perm8)
      let sgnC : Int := if inversionCount imgC % 2 == 0 then 1 else -1
      (tp + hs * sgnC, tm - hs * sgnC)
    else
      (tp, tm)
  ) (0, 0)

def wd4_chiral_trace_list : List (Int × Int) :=
  weyl_d4_elements.map (fun (p, s) => wd4ChiralTraces p s)

def wd4_norm_plus_sq : Int :=
  wd4_chiral_trace_list.foldl (fun acc (tp, _) => acc + tp * tp) 0

def wd4_norm_minus_sq : Int :=
  wd4_chiral_trace_list.foldl (fun acc (_, tm) => acc + tm * tm) 0

def wd4_cross_product : Int :=
  wd4_chiral_trace_list.foldl (fun acc (tp, tm) => acc + tp * tm) 0

set_option maxRecDepth 4096 in
theorem weyl_d4_even_parity_and_chiral_orthogonality_eq :
  (wd4_total_elements, wd4_all_even_parity, wd4_norm_plus_sq, wd4_norm_minus_sq, wd4_cross_product) =
    (192, true, 192, 192, 0) := by rfl
\end{lstlisting}

\appendix
\section{Verbatim Kernel Build and Zero-Axiom Verification Receipts}

All Lean~4 theorems presented in this paper were verified under Lean~4.35.0 and Lake~5.0.0 via \texttt{lake clean \&\& lake build}. Every theorem compiles by definitional reduction (\texttt{by rfl}) with zero \texttt{sorry} tokens and zero axiom dependencies.

\begin{lstlisting}[language={},basicstyle=\ttfamily\scriptsize,caption={Raw \texttt{lake clean \&\& lake build} and \texttt{\#print axioms} output for Paper~3 theorems.}]
[2/5] Built ShiodaFormalConjectures (64s)
info: ShiodaFormalConjectures.lean:469:0: 'weyl_d4_transition_matrix_invertible_cert' does not depend on any axioms
[3/5] Built HodgeVerification (22s)
info: HodgeVerification.lean:308:0: 'weil_w1_orthogonal_to_theta2' does not depend on any axioms
info: HodgeVerification.lean:309:0: 'weil_w2_orthogonal_to_theta2' does not depend on any axioms
info: HodgeVerification.lean:310:0: 'four_tuples_len' does not depend on any axioms
info: HodgeVerification.lean:311:0: 'zero_tuples_len' does not depend on any axioms
info: HodgeVerification.lean:325:0: 'clifford_even_dim_eq' does not depend on any axioms
info: HodgeVerification.lean:326:0: 'd4_triality_zero_triples_len' does not depend on any axioms
info: HodgeVerification.lean:460:0: 'spinor_to_weil_su22_branching_eq' does not depend on any axioms
info: HodgeVerification.lean:550:0: 'weyl_d4_conjugate_weil_axes_eq' does not depend on any axioms
info: HodgeVerification.lean:887:0: 'weyl_d4_spinor_weight_decomposition_eq' does not depend on any axioms
info: HodgeVerification.lean:950:0: 'weyl_d4_even_parity_and_chiral_orthogonality_eq' does not depend on any axioms
[4/5] Built LeanDemoProject (523ms)
Build completed successfully (5 jobs).
\end{lstlisting}

\begin{thebibliography}{9}
\bibitem{Deligne2000}
P.~Deligne, \emph{The Hodge Conjecture}, Clay Mathematics Institute Millennium Prize Problems, 2000.

\bibitem{Mumford1969}
D.~Mumford, \emph{A note of Shimura's paper ``Discontinuous groups and abelian varieties''}, Math. Ann. \textbf{181} (1969), 345--351.

\bibitem{vanGeemen1994}
B.~van Geemen, \emph{An introduction to the Hodge conjecture for abelian varieties}, in \emph{Algebraic Cycles and Hodge Theory}, Lecture Notes in Math. \textbf{1594}, Springer, 1994, 233--252.

\bibitem{Voisin2002}
C.~Voisin, \emph{Hodge Theory and Complex Algebraic Geometry} (Vols.~1--2), Cambridge Studies in Advanced Mathematics \textbf{76--77}, Cambridge University Press, 2002--2003.

\bibitem{Weil1977}
A.~Weil, \emph{Abelian varieties and the Hodge ring}, in \emph{Scientific Works / Collected Papers}, Vol.~3, Springer, 1979, 421--429.

\bibitem{Zarhin1984}
Yu.~G.~Zarhin, \emph{Weights of simple Lie algebras in the cohomology of algebraic varieties}, Izv. Akad. Nauk SSSR Ser. Mat. \textbf{48} (1984), 264--304.
\end{thebibliography}

\end{document}
